\input{../tex-inputs/latex-leseansicht-vorspann}

\section[Emil Marriot an Arthur Schnitzler, 27. 4. 1902]{L04332 Emil Marriot an Arthur Schnitzler, 27. 4. 1902}
\nopagebreak\mylabel{L04332v}
\rehead{ }\normalsize\beginnumbering\briefempfaengerindex{Schnitzler, Arthur@\textsc{Schnitzler, Arthur}!zzzMarriot, Emil@\emph{von Emil Marriot}!1902-04-272@{27. 4. 1902}|(be}
\toendnotes[C]{\smallbreak\pagebreak[2]}
\correspDesc{Versand  durch Emil Marriot am 27. 4. 1902 in Wien
\newline{}Erhalt  durch Arthur Schnitzler im Zeitraum [27. 4. 1902
                  – 30. 4. 1902?] in Wien}\toendnotes[C]{\smallbreak}
\Standort{DLA, A:Schnitzler, HS.1985.1.4008.}
\physDesc{Brief, 1 Blatt, 2 Seiten, 1508 Zeichen
\newline{}Handschrift: schwarze Tinte, lateinische Kurrent
\newline{}Schnitzler: 1) mit Bleistift Vermerk: »\textsc{Marriot}«  2) mit rotem Buntstift zwei Unterstreichungen}\toendnotes[C]{\smallbreak}
\pstart
           \raggedleft{}{\pb}Wien\oindex{Wien@\textbf{Wien}, \emph{Verwaltungsgebiet}|pw}, 27 IV. 1902.\pend
           
\pstart
           \raggedleft{}II/2 Schüttelstrasse 31 I.\oindex{Wien@\textbf{Wien}!II., Leopoldstadt@\textbf{II., Leopoldstadt}!Schüttelstraße 31@\textbf{Schüttelstraße 31}, \emph{Straße}|pw}\pend
           
\pstart{}Sehr geehrter Herr!\pend\vspace{0.5em}
\pstart
           Haben Sie sich auch schon die Frage vorgelegt, ob man einem Augenblicksimpuls folgen
               soll oder nicht? Mir scheint man soll’s and namentlich dann, wenn dabei nichts
               Schlimmes herauskommen kann. Dass Ihre »\label{K_L04332-1v}\edtext{lebendigen Stunden\pwindex{Schnitzler, Arthur 15. 5. 1862 Wien – 21. 10. 1931 ebd.@\textsc{Schnitzler, Arthur} (15. 5. 1862 Wien – 21. 10. 1931 ebd.), \emph{Schriftsteller, Mediziner}!Lebendige Stunden. Vier Einakter@\strich\emph{Lebendige Stunden. Vier Einakter}|pw}«, die ich in Berlin\oindex{Berlin@\textbf{Berlin}, \emph{Hauptstadt}|pw} sah}{\lemma{\textnormal{\emph{lebendigen … sah}}}\Cendnote{\textnormal{Die Uraufführung von \emph{Lebendige Stunden}\pwindex{Schnitzler, Arthur 15. 5. 1862 Wien – 21. 10. 1931 ebd.@\textsc{Schnitzler, Arthur} (15. 5. 1862 Wien – 21. 10. 1931 ebd.), \emph{Schriftsteller, Mediziner}!Lebendige Stunden. Vier Einakter@\strich\emph{Lebendige Stunden. Vier Einakter}|pwk}\eventindex{Deutsches Theater Berlin@\textbf{Deutsches Theater Berlin}!Uraufführung von Lebendige Stunden, 4.1.1902@Uraufführung von Lebendige Stunden, 4.1.1902|pwk} hatte am 4. 1. 1902 im Deutschen Theater in Berlin\oindex{Deutsches Theater Berlin@\textbf{Deutsches Theater Berlin}, \emph{Theater}|pwk} stattgefunden. }}}\label{K_L04332-1}, stark auf
               mich gewirkt haben, würde mich an sich noch nicht bestimmen, an Sie zu schreiben.
               Aber wir haben einen gemeinsamen Bekannten in Berlin\oindex{Berlin@\textbf{Berlin}, \emph{Hauptstadt}|pw}, Maximilian Harden\pwindex{Harden, Maximilian 20.\,10.\,1861 Berlin – 30.\,10.\,1927 Montana@\textsc{Harden, Maximilian} (20.\,10.\,1861 Berlin – 30.\,10.\,1927 Montana), \emph{Schriftsteller, Publizist}|pw}; und er
               fragte mich neulich in einem Brief, ob ich Sie kenne. Und da that es mir plötzlich
               leid, dass ich Ihnen niemals noch begegnet bin, auch vereinsamt und zurückgezogen wie
               ich lebe, keine Aussicht habe, Ihnen zufällig zu begegnen. Ich weiss nicht einmal, ob
               Sie sich gegenwärtig in Wien\oindex{Wien@\textbf{Wien}, \emph{Verwaltungsgebiet}|pw} aufhalten; ich lebe
               so gänzlich ausserhalb der literarischen Welt (Literaten sind mir im Allgemeinen
               unangenehm und Ihnen wohl auch), dass ich beinah von nichts und Niemanden etwas
               weiss. Doch \uline{Ihr} ernstes Schaffen, dass so wohlthuend
               vom Treiben unserer Reklamehelden sich abhebt, in jeder Beziehung abhebt, flösst mir
                  {\pb}auch für Ihre Persönlichkeit ein begreifliches
               Interesse ein. Es steht Ihnen ja selbstverständlich volkommen zu, diese Zeilen zu
               beachten oder zu ignoriren. Ich bin weder empfindlich noch verwöhnt. Aber dass ich
               Ihnen als Kollegin im Geist die Hand drücke and Ihnen für die schönen Stunden, die
               ich Ihren Werken verdanke, diesen Dank auch ausspreche, – Das darf ich doch, nicht
               wahr?\pend
           
\pstart
           Mit aller Hochachtung{\\[\baselineskip]}\spacefill\mbox{Emilie Mataja}{\\[\baselineskip]}\spacefill\mbox{(Emil Marriot.)}\pend
           \leftskip=0em{}\selectlanguage{ngerman}\endnumbering\briefempfaengerindex{Schnitzler, Arthur@\textsc{Schnitzler, Arthur}!zzzMarriot, Emil@\emph{von Emil Marriot}!1902-04-272@{27. 4. 1902}|)be}\mylabel{L04332h}
\begin{anhang}
\end{anhang}\newcommand{\dateiname}{L04332}\newcommand{\titel}{Emil Marriot an Arthur Schnitzler, 27. 4. 1902}\newcommand{\editorInnen}{Herausgegeben von Jahnke, SelmaMüller, Martin Anton}\input{../tex-inputs/latex-leseansicht-abspann}
